\begin{question}[index=17,points=10]
设 $f(x)$ 在 $(0,+\infty)$ 内有定义，$f'(1)=1$，对任意的正数 $x,y$ 满足 $f(xy)=\dfrac{f(x)}y+\dfrac{f(y)}x$.
\par（Ⅰ）求 $f(x)$；
\par（Ⅱ）求 $y=f(x)$ 的极值与拐点.
\end{question}
\QuestionBreak\WorkSpace

\begin{question}[index=18,points=12]
设 $f(x,y)=\begin{cases}\left(x^2+y^2+a\sqrt{|x|}+b\right)\dfrac{\sin(xy^2)}{x^2+y^4},&(x,y)\ne(0,0),\\0,&(x,y)=(0,0),\end{cases}$ 在点 $(0,0)$ 处可微，求 $a,b$ 的值.
\end{question}
\QuestionBreak\WorkSpace\WorkPageBreak

\begin{question}[index=19,points=12]
设 $f(x)$ 在 $[1,4]$ 上有连续导数，$f(4)-f(1)=1$，闭曲线 $L$ 是位于第一象限且由曲线 $y=x,y=4x,xy=1,xy=4$ 所围区域 $D$ 的正向边界，计算 $\displaystyle I=\oint_L\frac{f(xy)}y\dd y-\frac{f(xy)}x\dd x$.
\end{question}
\QuestionBreak\WorkSpace

\begin{question}[index=20,points=12]
设 $f(x)$ 在 $[-1,1]$ 上有连续的二阶导数，证明：
\par（Ⅰ）存在 $\xi\in[-1,1]$，使得 $\displaystyle\int_{-1}^1xf(x)\dd x=\frac13[2f'(\xi)+\xi f''(\xi)]$；
\par（Ⅱ）若 $f(x)$ 在 $(-1,1)$ 内取得极值，则存在一点 $\eta\in(-1,1)$，使得 $|2f'(\eta)+\eta f''(\eta)|\ge\dfrac12|f(1)-f(-1)|$.
\end{question}
\QuestionBreak\WorkSpace\WorkPageBreak

\begin{question}[index=21,points=12]
设矩阵 $\vect A=\mat{0&1&0\\1&0&2\\1&-1&2}$.
\par（Ⅰ）求 $3$ 阶可逆矩阵 $\vect P$ 与 $\vect Q$，使得 $\vect P\vect A\vect Q=\mat{1&0&0\\0&1&0\\0&0&0}$；
\par（Ⅱ）在第（Ⅰ）问的基础上，若记 $\vect B=(\vect Q\vect P)^{-1}$，且 $\vect A=\vect B\vect C$，求一个可逆矩阵 $\vect T$ 及对角矩阵 $\vect\Lambda$，使得 $\vect T^{-1}(\vect C+\vect E)^2\vect T=\vect\Lambda$.
\end{question}
\QuestionBreak\WorkSpace

\begin{question}[index=22,points=12]
设总体 $X\dist U[0,\theta]$，$X_1,X_2,\ldots,X_n$ 为来自总体 $X$ 的简单随机样本.
\par（Ⅰ）求 $\theta$ 的最大似然估计量 $\widehat\theta$，并说明 $\widehat\theta$ 是否为 $\theta$ 的无偏估计量；
\par（Ⅱ）记 $X_{(n)}=\max(X_1,X_2,\ldots,X_n)$，若选取 $T=\dfrac{X_{(n)}}\theta$ 为统计量，判断 $\theta$ 的 $1-\alpha$（$0<\alpha<1$）的置信区间是否为 $\displaystyle\left(\frac{X_{(n)}}{\sqrt[n]{1-\alpha/2}},\frac{X_{(n)}}{\sqrt[n]{\alpha/2}}\right)$.
\end{question}
\WorkSpace
